%! Permutations and Transposes — Unit 2, Lesson 2.4 — Warm-Up
\documentclass[10pt]{article}
\usepackage{linalg-article}
\usepackage{linalg-boxes}
\newcommand{\TallMath}[1]{$\displaystyle #1\rule[-1.4em]{0pt}{3.2em}$}
\newcommand{\xx}{\mathbf{x}}
\newcommand{\bb}{\mathbf{b}}
\newcommand{\tp}{\mathsf{T}}

\begin{document}

\pageheader{Unit 2, Lesson 2.4}{Warm-Up}
\namedateperiod

\vspace{0.12in}

\begin{notesbox}{Getting Ready: Skills We'll Use Today}
\small Today we handle the one case Lesson 2.3 skipped --- a $0$ in the pivot --- by \emph{swapping rows},
and we meet the \emph{transpose}. These three quick items are the pieces.

\vspace{4pt}
\textbf{1. A swap matrix} (Lesson 1.3). Multiply, one entry at a time:
\[
  \begin{bmatrix} 0 & 1 \\ 1 & 0 \end{bmatrix}
  \begin{bmatrix} 4 \\ 9 \end{bmatrix}
  = \begin{bmatrix}\rule{0pt}{1.1em}\blank{1.1cm}\\[2pt]\blank{1.1cm}\end{bmatrix}.
\]
Notice this matrix simply \emph{swaps} the two entries.

\vspace{4pt}
\textbf{2. A stuck pivot} (Lesson 2.1/2.3). To start elimination on
$A=\begin{bsmallmatrix} 0&3\\2&5 \end{bsmallmatrix}$, the first pivot is the top-left entry. That entry is
$\blank{1.1cm}$, and elimination cannot start because we would have to divide by
\rule[-2pt]{2.2cm}{0.6pt}.

\vspace{4pt}
\textbf{3. Reading entries by address} (Lesson 1.3). For $A=\begin{bsmallmatrix} 6&1\\4&9 \end{bsmallmatrix}$,
the entry in \emph{row 1, column 2} is $\blank{1.0cm}$ and the entry in \emph{row 2, column 1} is
$\blank{1.0cm}$. (Today the transpose will trade these two spots.)
\end{notesbox}

\end{document}
